\section{实验}

本节通过系统数值实验依次回答六个核心问题：第一，八年东海卫星遥感大数据能为航行器提供何种先验信息，深度学习时空外推能否替代现场采样（第5.2节）；第二，规划器在搜索用海洋上的自评得分随推演预算增加会产生多大的乐观偏差，留出海洋能否给出无偏评判（第5.3节）；第三，在四种机载硬件与两种真值模型下，出发前预演与需要历史任务的方法、在线元推理及启发式巡测规则相比表现如何（第5.4节）；第四，经典单步元推理为何失效，留出估值与出发前预演通过何种机制实现硬件自适应（第5.5节）；第五，预演中的假想真值数量、不确定性标定与评估尺度各起多大作用（第5.6节）；第六，在算力速度、底图偏离度与底图新鲜度构成的二维参数空间中，深想与停下复测的适用边界何在（第5.7节）。

\subsection{实验设置与对比方法}

全部实验均基于2018至2026年东海3197天卫星叶绿素与海表温度数据集构建的422个赤潮任务窗口开展，其中2018至2022年的274个开发集窗口用于统计底图误差、绘制算力相图及训练历史基线，2023至2026年的148个测试集窗口用于主实验与消融检验（表\ref{tab:settings}）。实验涵盖四种机载硬件配置：慢计算机（低功耗处理器每分钟推演5次，高功耗芯片$+20\text{ W}$）、中等计算机（20次/分，$+20\text{ W}$）、快计算机（100次/分，$+20\text{ W}$）以及省电加速器（20次/分，$+5\text{ W}$），并在“当天卫星图 + 小尺度变化”（主设置，$s_{\text{sat}}=0.15$）与“当天卫星图”（全实测数据设置，$s_{\text{sat}}=0$）两种真值模型下开展评估。为保证严格配对比较，同一任务窗口上的所有控制方式均共享完全相同的真实海洋、入水起点与荧光计随机噪声序列，且每次任务的随机数种子固定；控制器间时间平均遗憾配对差值的95\%置信区间通过按赤潮周整组重抽样的自助法计算（同一周内的任务窗口一并重抽样以计入周内空间相关性）。

\begin{table}[htbp]
  \centering
  \caption{东海赤潮应急巡测实验的核心参数与评价设置。}
  \label{tab:settings}
  \begin{tabular}{@{}p{0.17\linewidth}p{0.79\linewidth}@{}}
\toprule
项目 & 设置 \\
\midrule
卫星数据 & 美国国家海洋和大气管理局海岸观测网多星融合叶绿素 a 产品（经验正交函数插补去云），2 km、逐日，2018-01-01 至 2026-10-23 共 3197 天；喷气推进实验室多尺度融合海表温度，0.01° \\
任务窗口 & 每年 4--6 月浙江外海陆架（26.5°--31.5°N，121.5°--125°E）局部峰值超过 5 mg/m$^3$ 处，80 km $\times$ 80 km（40 $\times$ 40 个 2 km 网格）；开发 274 个（2018--2022 年），测试 148 个（2023--2026 年） \\
真实海洋 & 任务当天卫星图叠加卫星分辨不出的小尺度变化（标准差 0.15，相关尺度 6 km，主设置）；或任务当天卫星图本身 \\
航行器 & 航速 1.5 m/s，电池 1 kWh；平台 15 W，推进 30 W，高功耗芯片 +20 W（省电加速器 +5 W）；任务 24 h，耗尽电量即停 \\
决策周期 & 20 min；走一段约 1.8 km、采样约 2 次；停下漂流采样 5 次；荧光计单次误差 0.1 \\
想象海洋 & 48 个（32 个用于搜索航线，16 个留出），由 3 天前卫星图与数据标定的误差生成，空间相关长度 10 km，随每个样本同化更新 \\
规划器 & 交叉熵搜索经两个航路点的折线航线，跨周期热启动；低功耗处理器每分钟 5、20 或 100 次推演（条件扫描 5--200 次），高功耗芯片为其 50 倍 \\
出发前预演 & 4 个留出海洋作假想真值 $\times$ 5 种用时策略，各完整模拟 24 h，选平均假想遗憾最小者 \\
评价 & 时间平均遗憾（对数浓度，越低越好）；配对差值的 95\% 置信区间由按赤潮周整组的自助重抽样得到 \\
\bottomrule
\end{tabular}

\end{table}

实验将“出发前预演（本文）”与五类控制方式进行全面对比。第一类为五种固定时间分配策略，即“一直走”、“一直深想”、“三段深想一段”、“三段停下深想一段”和“三段停下复测一段”，以及不经预演的“五种策略随机选”和“三种不停策略随机选”。第二类为需要该型航行器历史巡测任务的方法，包括“调好的固定策略”（在2018至2022年120个开发集任务上为每种硬件和真值模型选出平均遗憾最低的固定策略，再直接用于测试集，属于强基线）和“按历史任务学习的选择器”（参照学习型元推理思想\cite{magnino2026metareasoning}，在同样120个开发集任务上训练梯度提升回归模型，利用出发前已知的8项旧图特征——最大值、中位数、极差、标准差、起点至旧图峰值距离、可航行网格占比、起点处旧图浓度及接近峰值网格占比——预测五种策略的遗憾并为每次任务挑选预测遗憾最低者）。第三类为在线元推理，包括经典“单步元推理”\cite{russell1991right}与本文改进的“单步元推理（留出估值）”。第四类为不开展多步航线规划的经典海洋巡测规则（航行器全程保持巡航采样并以相同局地化集合滤波更新地图），包括“割草机式全覆盖”（从入水边缘起按间距$16\text{ km}$平行往返扫测）\cite{galceran2013coverage}、“贪心：驶向估计最浓处”（每个周期驶向集合均值最高的未访网格）、“上置信界采样”（驶向集合均值加2倍标准差最高的未访网格）\cite{srinivas2012gpucb}以及“汤普森采样”（每个周期随机抽取一个想象海洋并驶向其最浓未访网格）\cite{russo2018thompson}。第五类为仅作性能下界参考的“每次任务事后最优（参考）”，即在每次测试任务结束后事后挑出五种固定策略中实际遗憾最低者。

\subsection{海洋大数据对底图不确定性的标定作用与预测边界}

八年卫星叶绿素大数据能够精确标定滞后底图在不同季节与滞后天数下的误差分布，但即便是包含550万参数的深度时空预测模型也无法改善1至3天的底图外推精度，小尺度藻华结构必须依靠现场采样获取（图\ref{fig:data}）。在全部3197天数据上统计相隔3天的对数叶绿素变化标准差可见，东海表层叶绿素的3天演变幅度呈现出清晰的季节与空间分异：全海域中位数在赤潮高发的春季最高（0.127），夏季（0.107）与秋季（0.103）次之，冬季降至0.081，而在夏季浙闽沿岸与长江冲淡水锋面一带形成了3天变化标准差高达0.25至0.3的强变率条带（图\ref{fig:data}a至图\ref{fig:data}d）。在274个春季开发集窗口内，任务当天卫星图减去1、2、3天前旧图的误差标准差随滞后天数单调递增，依次为$0.125$、$0.176$与$0.210$（均值偏差依次为$+0.016$、$+0.032$与$+0.045$，空间相关尺度约为5个网格即$10\text{ km}$），且同日卫星图与现场海水之间还存在约0.15的对数浓度差异\cite{ye2021cocts}（图\ref{fig:data}e）。

\begin{figure}[htbp]
  \centering
  \includegraphics[width=\linewidth]{fig_data}
  \caption{八年东海卫星叶绿素大数据提供的先验误差标定及其外推预测边界。（a）至（d）2018至2026年全部3197天对数叶绿素3天变化标准差的四季空间分布（中位数春季0.127、夏季0.107、秋季0.103、冬季0.081），春季子图标出了274个开发集窗口与148个测试集窗口的位置。（e）开发集窗口内滞后1、2、3天旧卫星图的误差标准差（0.125、0.176、0.210）及同日卫星图相对于现场实测水体的约0.15偏差。（f）2023至2026年测试集上深度时空预测模型与直接沿用最新卫星图在提前1、2、3天预报上的均方根误差对比（0.1229对0.1226、0.1439对0.1439、0.1487对0.1474）。数据能够可靠标定旧图的不确定性幅值，但无法替代现场观测。}
  \label{fig:data}
\end{figure}

为检验能否在航行器出发前用深度学习模型将3天前旧图向前推演以缩减滞后误差，本文训练了一个550万参数的联合嵌入预测架构模型\cite{assran2025vjepa2}，以过去3天的叶绿素和海表温度异常场为输入，使用2018至2021年3月至6月数据训练并在2022年上验证（验证集误差仅比沿用最新底图低1.4\%）。然而在2023至2026年测试年份上，该模型提前1、2、3天预报的均方根误差分别为$0.1229$、$0.1439$与$0.1487$（集合离散度与误差之比为$0.70\text{--}0.93$），与直接沿用最新可用卫星图的持久性基线（$0.1226$、$0.1439$与$0.1474$）完全持平（图\ref{fig:data}f）。这一结果从物理与数据层面确立了海洋大数据在本问题中的真正角色：由于云区插补产品平滑了日内潮汐混合与亚中尺度梯度，数值外推无法凭空恢复小尺度峰值，但多年遥感统计能够为生成想象海洋提供准确的先验不确定性标尺。

\subsection{规划器自评得分的乐观偏差与留出海洋的无偏评判}

随着单次决策推演预算的增加，规划器在搜索用海洋上的自评得分因优化器诅咒而持续虚高，而在未参与搜索的留出海洋上评判则能消除这一偏差并准确反映航线真实价值的饱和过程（图\ref{fig:curse}）。在80个开发集窗口（每窗口3个随机种子，携带3天前卫星图，评估任务前8小时，主真值模型）上，当每次决策的推演次数从64次增至2000次再扩展至200,000次（跨度达三个半数量级）时，所选航线在真实海洋中的实际价值从$0.749$升至$0.793$，随后进入平台期并平缓收敛至$0.803$（图\ref{fig:curse}a黑线）。相反，交叉熵规划器在32个搜索用海洋上的内部自评得分却从$0.762$升至$0.833$并一路攀升至$0.917$（图\ref{fig:curse}a红线），在最大预算下比真实航线价值高估达$0.114$。当把同一批搜索选出的航线放到16个留出海洋上评估时，留出评判均值依次为$0.745$、$0.794$与$0.802$（图\ref{fig:curse}a蓝虚线），在全部推演预算下与真实价值之差均不超过$0.005$（图\ref{fig:curse}b）。

\begin{figure}[htbp]
  \centering
  \includegraphics[width=\linewidth]{fig_curse}
  \caption{推演预算从64次扩展至200,000次时规划器自评得分的乐观偏差与留出海洋的无偏评判（80个开发集窗口$\times 3$个种子，3天前卫星图，任务前8小时）。（a）所选航线的时间平均最佳对数叶绿素价值：搜索用海洋上的自评得分从0.762持续升至0.917，而真实航线价值（0.749升至0.803）与留出海洋评判（0.745升至0.802）在约2000次推演后一同趋于饱和。（b）估计价值减去真实价值的偏差曲线：使用数据标定不确定性时留出评判偏差保持在0.005以内，将不确定性设小一半或放大至1.5倍虽产生$-0.06$与$+0.07$的整体平移，但仍能准确捕捉增加算力的相对增益。}
  \label{fig:curse}
\end{figure}

进一步改变先验标准差幅值的实验表明，即使不确定性标定存在一定偏差，留出海洋对增加算力所获边际收益的判断依然稳健（图\ref{fig:curse}b）。当人为将先验标准差$\sigma_3$减半或放大至1.5倍时，留出评判绝对值虽分别产生约$-0.06$与$+0.07$的系统性平移，但推演预算从64次增至200,000次时留出海洋估计的价值增量（分别为$+0.059$与$+0.054$）依然与真实海洋中的实际价值增量（$+0.054$与$+0.047$）高度吻合。这证实将搜索用海洋与留出海洋分离能够彻底切断优化器对样本噪声的迎合链条，为衡量计算价值提供可靠的独立标尺。

\subsection{多硬件平台与双真值模型下的主实验对比}

在2023至2026年全部148个测试任务上，出发前预演在四种机载硬件与两种真值模型下均无需任何历史巡测任务即达到“调好的固定策略”与“按历史任务学习的选择器”的同等水平，并一致显著优于经典“单步元推理”和各类无规划巡测规则（图\ref{fig:main}，表\ref{tab:main}）。在“当天卫星图 + 小尺度变化”主设置下，慢计算机、中等计算机、快计算机与省电加速器上“出发前预演（本文）”的时间平均遗憾分别为$0.276$、$0.265$、$0.268$与$0.266$。与各硬件在120个历史开发任务上专门选出的“调好的固定策略”（依次为“三段深想一段”$0.270$、“三段深想一段”$0.269$、“一直走”$0.267$与“一直深想”$0.267$）相比，出发前预演减去调好策略的配对差值依次为$+0.007$（95\%置信区间$[-0.000, +0.014]$）、$-0.003$（$[-0.009, +0.003]$）、$+0.002$（$[-0.005, +0.008]$）与$-0.001$（$[-0.006, +0.005]$），与“按历史任务学习的选择器”（遗憾依次为$0.278$、$0.269$、$0.273$与$0.266$）的配对差值依次为$-0.001$、$-0.004$、$-0.004$与$+0.000$，全部八组对比的95\%置信区间均包含0。

\begin{figure}[htbp]
  \centering
  \includegraphics[width=\linewidth]{fig_main}
  \caption{148个测试任务上各对比控制方式减去“出发前预演（本文）”的时间平均遗憾配对差值及95\%置信区间（正值表示差于本文方法）。四列对应慢计算机、中等计算机、快计算机与省电加速器，实心圆点为“当天卫星图 + 小尺度变化”主设置，空心圆点为“当天卫星图”真值设置。出发前预演在所有硬件与两种真值模型下均匹配需要历史任务的两种方法，并显著优于经典单步元推理及所有巡测规则（仅慢计算机主设置下与贪心规则持平）。}
  \label{fig:main}
\end{figure}

\begin{table}[htbp]
  \centering
  \caption{2023至2026年148个测试任务上四种机载硬件与两种真值模型下各控制方式的时间平均遗憾（越低越好，每列最优可运行控制器加粗，本文方法行加底纹）。}
  \label{tab:main}
  \small\resizebox{\linewidth}{!}{\begin{tabular}{@{}llcccccccc@{}}
\toprule
 & & \multicolumn{4}{c}{真值：当天卫星图 + 小尺度变化} & \multicolumn{4}{c}{真值：当天卫星图} \\
\cmidrule(lr){3-6}\cmidrule(l){7-10}
类别 & 控制方式 & 慢计算机 & 中等计算机 & 快计算机 & 省电加速器 & 慢计算机 & 中等计算机 & 快计算机 & 省电加速器 \\
\midrule
固定策略 & 一直走 & 0.374 & 0.267 & 0.267 & 0.267 & 0.281 & 0.196 & 0.194 & 0.196 \\
 & 一直深想 & 0.273 & 0.269 & 0.269 & 0.267 & 0.195 & 0.194 & 0.193 & \textbf{0.189} \\
 & 三段深想一段 & \textbf{0.270} & 0.269 & \textbf{0.265} & 0.266 & \textbf{0.194} & \textbf{0.190} & \textbf{0.191} & 0.191 \\
 & 三段停下深想一段 & 0.331 & 0.328 & 0.327 & 0.327 & 0.252 & 0.248 & 0.245 & 0.247 \\
 & 三段停下复测一段 & 0.408 & 0.335 & 0.331 & 0.335 & 0.319 & 0.251 & 0.248 & 0.251 \\
\midrule
需要历史任务 & 调好的固定策略$^{a}$ & \textbf{0.270} & 0.269 & 0.267 & 0.267 & \textbf{0.194} & 0.196 & \textbf{0.191} & 0.191 \\
 & 按历史任务学习的选择器 & 0.278 & 0.269 & 0.273 & \textbf{0.266} & 0.200 & 0.197 & \textbf{0.191} & 0.189 \\
\midrule
在线元推理 & 单步元推理 & 0.287 & 0.287 & 0.282 & 0.278 & 0.202 & 0.202 & 0.204 & 0.199 \\
 & 单步元推理（留出估值） & 0.278 & 0.270 & 0.270 & 0.267 & -- & -- & -- & -- \\
\midrule
巡测规则 & 割草机式全覆盖 & 0.424 & 0.424 & 0.424 & 0.424 & 0.328 & 0.328 & 0.328 & 0.328 \\
 & 贪心：驶向估计最浓处 & 0.285 & 0.285 & 0.285 & 0.285 & 0.206 & 0.206 & 0.206 & 0.206 \\
 & 上置信界采样 & 0.299 & 0.299 & 0.299 & 0.299 & 0.224 & 0.224 & 0.224 & 0.224 \\
 & 汤普森采样 & 0.332 & 0.332 & 0.332 & 0.332 & 0.256 & 0.256 & 0.256 & 0.256 \\
\midrule
\rowcolor{violet!8}本文 & 出发前预演 & 0.276 & \textbf{0.265} & 0.268 & \textbf{0.266} & 0.197 & 0.191 & 0.195 & 0.190 \\
\midrule
参考 & 每次任务事后最优 & 0.250 & 0.243 & 0.243 & 0.242 & 0.183 & 0.179 & 0.180 & 0.179 \\
\bottomrule
\end{tabular}
}
  \par\vspace{2pt}\footnotesize 注：$^a$ 调好的固定策略在主设置下依次为三段深想一段、三段深想一段、一直走、一直深想，在“当天卫星图”真值下依次为三段深想一段、一直走、三段深想一段、三段深想一段。
\end{table}

与在线元推理和启发式巡测规则相比，出发前预演与多步航线规划的优势在统计上高度显著。在主设置下，经典“单步元推理”在四种硬件上的遗憾高达$0.287$、$0.287$、$0.282$与$0.278$，出发前预演相对其降低遗憾$-0.011$（$[-0.019, -0.002]$）、$-0.022$（$[-0.030, -0.013]$）、$-0.013$（$[-0.020, -0.007]$）与$-0.012$（$[-0.019, -0.004]$），95\%置信区间全部严格位于0以下。四种不含多步路线规划的巡测规则因始终只做单步或固定航线机动，在四种硬件上表现相同：“贪心：驶向估计最浓处”遗憾为$0.285$，“上置信界采样”为$0.299$，“汤普森采样”为$0.332$，“割草机式全覆盖”高达$0.424$；出发前预演除在慢计算机上与贪心规则差值为$-0.009$（$[-0.022, +0.004]$）外，在其余三种硬件上比贪心规则低$-0.017$至$-0.020$，并在所有硬件上比上置信界采样、汤普森采样和割草机式全覆盖分别低$-0.023$至$-0.034$、$-0.056$至$-0.067$以及$-0.148$至$-0.159$（置信区间全部低于0）。

在直接以任务当日实测卫星图为真值的“当天卫星图”设置下，由于不存在卫星未分辨的亚网格随机场，所有控制器的绝对遗憾均整体下降，但各方法之间的优劣排序完全不变（表\ref{tab:main}右半部）。“出发前预演（本文）”在四种硬件上取得$0.197$、$0.191$、$0.195$与$0.191$，与“调好的固定策略”（$0.194$、$0.196$、$0.191$与$0.191$）之差为$-0.004$至$+0.003$，与“按历史任务学习的选择器”之差为$-0.005$至$+0.003$（置信区间均包含0），同时比“单步元推理”（$0.202$、$0.202$、$0.204$与$0.199$）低$-0.006$至$-0.010$，比四种巡测规则低$-0.010$至$-0.139$（置信区间全部低于0）。此外，两种真值模型下148个测试任务之间的遗憾标准差分别高达约$0.13$与$0.11\text{--}0.12$，而三种不停航策略之间的平均差距（除慢计算机外）不超过$0.01$，这直接解释了为何“按历史任务学习的选择器”虽利用了8项旧图特征和120个历史训练任务，却依然无法在测试集上超越单一的调好策略。同时，对慢计算机全部148个测试任务重新独立运行预演，所得每次任务的遗憾数值与首次运行完全一致，证实了实验结果的精确可重复性。

\subsection{作用机制分析：在线过度规划的成因与预演的硬件自适应}

经典单步元推理之所以表现不佳，是因为轻信搜索集自评得分使其在80\%至97\%的周期内盲目开启高功耗芯片并在约三分之一的周期内停下复测，而改用留出估值或采用出发前预演则能大幅压低无效计算与停船，并使所选策略随硬件算力和功耗自适应调整（图\ref{fig:mech}）。在主设置下，“单步元推理”在慢、中、快计算机与省电加速器上开启高功耗芯片的周期占比分别高达80\%、93\%、97\%与94%（图\ref{fig:mech}a），处于停下漂流状态的周期占比均达33\%至34\%（图\ref{fig:mech}b）。当保持单步元推理其余流程完全不变、仅将计算收益与复测收益改在16个留出海洋上评估后，“单步元推理（留出估值）”开启高功耗芯片的比例骤降至30\%、38\%、44\%与39\%，停下比例减半至15\%至18\%，其时间平均遗憾在四种硬件上依次降至$0.278$、$0.270$、$0.269$与$0.267$（比经典单步元推理分别降低$0.009$、$0.017$、$0.013$与$0.011$），与“出发前预演（本文）”的差距缩小至$0.001$至$0.005$以内（所有95\%置信区间均包含0）。这一对照实验直接证明，经典元推理的过度规划正是由搜索样本上的自评乐观偏差引起的，而剩余15\%至18\%的停下周期则是其略逊于预演（预演停下比例仅$0.5\%\text{--}1.8\%$）的唯一残余损失来源。

\begin{figure}[htbp]
  \centering
  \includegraphics[width=\linewidth]{fig_mech}
  \caption{在线元推理过度规划的成因与出发前预演的硬件自适应机制（主真值设置，148个测试任务）。（a）高功耗芯片开启周期占比与（b）停下复测周期占比：经典单步元推理在80\%至97\%的周期内开启高功耗芯片并在33\%至34\%的周期内停下，而单步元推理（留出估值）与出发前预演大幅削减了无效深想与停船。（c）出发前预演在四种硬件上选出各候选策略的任务占比：随基础算力提升，“一直走”占比从1\%升至30\%再升至47\%，而在省电加速器上“一直深想”占比从25\%升至43\%。}
  \label{fig:mech}
\end{figure}

从出发前预演在148个测试任务上选出的策略分布（图\ref{fig:mech}c）可以清晰看出其适配不同硬件的物理机制。在后台算力极慢（5次/分）、“一直走”遗憾高达$0.374$（比调好策略差$+0.104$）的慢计算机上，预演选出“三段深想一段”占55\%、“一直深想”占39\%，“一直走”仅占1\%，停下类策略仅占4\%。随着低功耗处理器速度提升至中等计算机（20次/分）和快计算机（100次/分），后台规划已接近航线价值饱和区，预演选出“一直走”的比例随之从1\%跃升至30\%（“三段深想一段”43\%、“一直深想”25\%、停下类1\%）并进一步升至快计算机上的47\%（“三段深想一段”28\%、“一直深想”22\%、停下类3\%）。而在算力与中等计算机相同（20次/分）、但高功耗芯片功耗从$+20\text{ W}$降至$+5\text{ W}$的省电加速器上，由于开启深想对续航的惩罚大幅减轻，预演选出“一直深想”的比例从中等计算机的25\%上升至43\%（“三段深想一段”35\%、“一直走”20\%、停下类1\%）。由此可见，预演的核心收益在于准确识别并剔除当前硬件下的劣质选项（所有硬件上的两种停下策略以及慢计算机上的“一直走”），从而显著优于“五种策略随机选”（遗憾$0.331$、$0.293$、$0.292$与$0.292$）和慢计算机上的“三种不停策略随机选”（$0.306$）。而在真实平均差距不超过$0.004$的其余三种硬件的不停策略之间，预演的单任务区分能力接近随机抽取：“出发前预演（本文）”与“预演（以搜索过的海洋为真值）”在单个任务上选出同一策略的比例仅为37\%至41\%，且两者平均遗憾仅相差$-0.004$至$+0.004$（置信区间均包含0），这是因为预演依据假想真值上的最终任务遗憾而非规划器自评得分来比较策略，本身已对规划器自评泡沫免疫。

\begin{figure}[htbp]
  \centering
  \includegraphics[width=\linewidth]{fig_example}
  \caption{一次典型测试任务（2023年4月15日，慢计算机）中“一直走”与“三段深想一段”的航迹及实时遗憾对比。（a）航行器出发时携带的3天前卫星图及东侧边缘入水起点（白色三角形）。（b）真实海洋叶绿素分布（西南角五角星为真实最浓处）以及“一直走”（蓝线）与“三段深想一段”（橙线，加粗段为开启高功耗芯片深想）的24小时航迹。（c）24小时内报告位置的实时遗憾曲线：“一直走”因搜索不足在窗口中部绕圈，遗憾停留在0.4附近（均值0.43），“三段深想一段”横穿窗口并在约19小时后抵达真实最浓处（均值降至0.31）。}
  \label{fig:example}
\end{figure}

一次具体测试任务（2023年4月15日，慢计算机，按日期顺序首个两种策略差距超过0.1的任务）直观展示了慢速处理器上开启高功耗芯片深想如何改变航行轨迹（图\ref{fig:example}）。该任务中航行器从窗口东侧边缘入水，3天前旧图显示西侧与西南侧浓度较高（图\ref{fig:example}a），而真实海洋的最浓核心区位于西南角（图\ref{fig:example}b五角星）。在慢计算机上，“一直走”每20分钟仅能完成100次后台推演，粗糙的航路点搜索使其航行至窗口中部次高值区后便陷入局部徘徊（图\ref{fig:example}b蓝线），24小时内报告误差始终停留在0.4左右，时间平均遗憾达$0.43$（图\ref{fig:example}c）。相比之下，“三段深想一段”每3个周期开启一次高功耗芯片完成5000次推演（图\ref{fig:example}b橙色加粗段），充裕的搜索预算使其及时识别出西南方向的远端高梯度潜力并果断横穿整个窗口，在约19小时后准确抵达西南角真实最浓处，将24小时时间平均遗憾压低至$0.31$。

\subsection{消融实验：假想真值数量、不确定性标定与评估尺度}

在慢计算机与中等计算机上开展的单因素消融实验表明，充足的假想真值数量、由卫星大数据标定的先验不确定性幅值以及全任务预演或留出估值机制，均是保障方法稳健性的必要环节（图\ref{fig:ablation}）。首先，当把预演使用的留出假想真值数量从默认的4个缩减为1个或2个时，慢计算机与中等计算机上的时间平均遗憾分别劣化$+0.013$（$[-0.000, +0.027]$）与$+0.010$（$[-0.001, +0.020]$），以及$+0.010$（$[-0.001, +0.022]$）与$+0.006$（$[-0.002, +0.013]$），说明单个假想真值的空间涨落较大，容易误选次优策略；而将假想真值增至8个时，慢计算机遗憾虽进一步微降$-0.007$（$[-0.012, -0.001]$），但中等计算机变化为$+0.003$（$[-0.004, +0.009]$）且计算耗时翻倍，因此4个假想真值是兼顾精度与岸基计算效率的合理选择。其次，将4个假想真值改为从32个搜索用海洋中抽取的“预演（以搜索过的海洋为真值）”，其遗憾与完整方法仅相差$-0.004$与$+0.003$（置信区间均包含0），再次印证只要以假想真值上的最终任务结果而非规划器自评得分作为评判标准，即可规避自评过拟合。

\begin{figure}[htbp]
  \centering
  \includegraphics[width=\linewidth]{fig_ablation}
  \caption{慢计算机与中等计算机（主真值设置，148个测试任务）上的消融实验结果（各变体减去完整“出发前预演（本文）”的时间平均遗憾配对差值及95\%置信区间）。减少假想真值数量、将先验不确定性错误放大至2倍或取消预演改为随机选择均导致性能劣化，而在任务尺度预演、在线周期级留出估值与需要历史任务的方法之间表现相当。}
  \label{fig:ablation}
\end{figure}

先验不确定性的遥感数据标定对同化与预演质量具有决定性影响。当不依赖数据标定而人为将先验标准差$\sigma_3$减半时，慢、中计算机遗憾变化为$+0.002$与$+0.004$（置信区间包含0），但当将先验标准差放大至标定值的2倍时，慢计算机与中等计算机的遗憾分别显著恶化$+0.011$（$[+0.001, +0.020]$）与$+0.020$（$[+0.011, +0.030]$）。在2倍不确定性下，过大的先验噪声淹没了旧卫星图的大尺度梯度信息，不仅使所有固定策略本身的表现全面退化（中等计算机上三种不停策略的遗憾从$0.267\text{--}0.269$恶化至$0.275\text{--}0.283$），而且使预演选出的策略（$0.286$）落后于最佳固定策略（$0.275$）。最后，将全任务出发前预演替换为在线逐周期的“单步元推理（留出估值）”时，慢、中计算机差值仅为$+0.002$与$+0.005$（置信区间包含0），与需要历史任务的“调好的固定策略”（$-0.007$与$+0.003$）和“按历史任务学习的选择器”（$+0.001$与$+0.004$）同样处于统计等价区间，而完全取消预演改为“五种策略随机选”则会带来$+0.055$与$+0.028$的显著损失（“三种不停策略随机选”在慢计算机上亦损失$+0.029$）。

\subsection{适用边界：算力、底图偏离度与底图新鲜度的影响}

在低功耗处理器速度与真实海洋偏离卫星图程度构成的二维参数空间中，开启高功耗芯片深想仅在基础算力低于航线价值饱和阈值时产生显著收益且底图越不准收益越大，而中途停下原位复测在所有参数组合下均造成显著损失（图\ref{fig:phase}）。在由低功耗处理器推演速度（5、10、20、50、100、200次/分）与真实海洋偏离同日卫星图程度（$s_{\text{sat}}$取0、0.05、0.15、0.30，每格60个开发集任务）组成的网格相图上，比较较好的深想策略（“一直深想”与“三段深想一段”中较优者）与“一直走”可见（图\ref{fig:phase}a）：唯有在最慢的5次/分处理器上（后台每20分钟仅推演100次，远低于图\ref{fig:curse}a中约2000次推演的饱和点），开启高功耗芯片深想才稳定显著优于“一直走”，且当偏离程度从0.05增至0.15再增至0.30时，深想降低的时间平均遗憾从$-0.086$扩大至$-0.109$再扩大至$-0.138$，表明底图越不可信，越需要充足的算力结合现场同化重塑航线。在10次/分处理器上，仅偏离0.15一格存在微弱收益（$-0.017$，$[-0.035, -0.001]$），而从20次/分（每周期后台推演400次）起直至200次/分，所有偏离程度下深想均不再带来显著收益。当真实海洋偏离程度为0（即“当天卫星图”本身作为真值）时，相图第四行的结果与上述规律一致：\TBD{待补偏离度为0行的具体数值}。

\begin{figure}[htbp]
  \centering
  \includegraphics[width=\linewidth]{fig_phase}
  \caption{低功耗处理器速度（每分钟5至200次推演）与真实海洋偏离卫星图程度（0、0.05、0.15、0.30）二维网格上的策略收益相图（每格60个开发集任务）。（a）最佳深想策略减去“一直走”的时间平均遗憾差值（加框格点表示95\%置信区间低于0，即深想显著有益）：深想仅在基础算力过慢（5次/分）时显著获益，且底图偏离越大收益越强。（b）较好停下策略减去最佳不停策略的遗憾差值：所有格点差值均为显著正值（$+0.063$至$+0.118$），表明停下原位复测在任何算力与底图误差下均得不偿失。}
  \label{fig:phase}
\end{figure}

相图右半部（图\ref{fig:phase}b）揭示了另一条普适规律：在全部算力速度与底图偏离度格点上，两种停下策略中较好者均比最佳不停策略高出$+0.063$至$+0.118$的时间平均遗憾，所有格点的95\%置信区间均严格大于0。这一现象可由不确定性空间尺度与时间积分目标共同解释：在原地停下漂流20分钟连续采样5次，虽能将该单点处的荧光计测量标准差从$0.1$压低至$0.1/\sqrt{5}\approx 0.045$，但制约航行器找到全窗最浓水体的瓶颈并非已访点的传感器噪声，而是广阔未访水域高达约$0.26$的先验不确定性，且将三分之一的作业时间用于原地漂流会严重推迟抵达核心浓水区的时刻，直接招致时间平均遗憾的沉重惩罚。此外，当航行器装入的是滞后仅1天的新鲜卫星图（$k=1$，$s_1=0.125$，主真值设置）时，在148个测试任务上各控制器绝对遗憾整体下降：在慢计算机上，“调好的固定策略”为“一直深想”（$0.250$），“出发前预演（本文）”取得$0.251$（与调好策略差值$+0.001$，$[-0.005, +0.011]$），依然优于“单步元推理”的$0.260$（差值$-0.009$，$[-0.018, +0.001]$），而在中等计算机上，\TBD{待补1天前卫星图在中等计算机上的具体数值}。

\subsection{局限性}

本研究在海洋环境表征、运动规划维度与预演执行机制上仍存在若干简化，界定了当前结论的适用范围。首先，实验中的真实海洋是由任务当天的真实卫星叶绿素图（叠加或不叠加按实测文献标定的小尺度随机场）构成的，并非现场走航实测剖面，且二维表层场与随表层水体漂移的参考系忽略了窗口内的水平流剪切以及东海原甲藻的昼夜垂直迁移行为，双航路点折线规划器也是对水下三维流场路径规划的简化替代。其次，出发前预演仅在下水前由母船或岸基计算机执行一次，未在24小时巡测中途根据同化更新后的地图重新切换策略，且4个假想真值的预演主要实现了针对机载硬件平台特性的策略适配，尚不足以在平均差距微小的几种不停策略之间针对单次任务做出优于随机的精细区分（主设置下距离“每次任务事后最优（参考）”仍有$0.022\text{--}0.027$的差距）。最后，预演内部的假想遗憾绝对值比真实海洋中的实际遗憾系统性偏高约$0.05\text{--}0.07$，但由于这一偏移对所有候选策略均等发生，并未改变各策略之间的相对优劣排序。
